% Einheit 1 von 4 – Zählen und Ergebnismengen (bank/wahrscheinlichkeit/e1.jsonl, zusammenbau v0.1)
%% TODO Zweigzeile Teil 1: was hier gelernt wird (Satz in Schülersprache aus dem Einheitstitel) – Einheit 1
\einheitenkopf[e1]{Einheit 1 von 4 · Zählen und Ergebnismengen}
\zweigzeile{ab Kl. 5, je nach Buch bis Kl. 10 (am Gymnasium ab Kl. 6) · P10}
\setcounter{aufgabe}{7}

%% TODO Titel: Ich-kann-Satz fehlt in der Bank (Platzhalter: Kettenname) – Nr. 8
%% TODO Anweisung: ein Satz über der ersten Teilaufgabe fehlt in der Bank – Nr. 8
\begin{aufgabe}{Zählen}
%% TODO \rechenplatz{n}: Schrittzahl des Grundfalls fehlt in der Bank (nur bei fester Schreibform, 2.3 b)
\begin{teile}
\teil Paul hat zwei Mützen (grau, grün) und zwei Schals (gelb, weiß). Schreibe alle Paare aus Mütze und Schal auf – erst alle mit der grauen Mütze.
\teil Lena hat 3 Hosen und 2 Pullover. Wie viele Kombinationen aus Hose und Pullover gibt es? \leerfeld Kombinationen
\teil Tom hat 5 Hosen und 7 T-Shirts. Wie viele Kombinationen gibt es?
\teil Anna, Ben und Cem stellen sich in eine Reihe. Schreibe alle Reihenfolgen auf. Wie viele sind es?
\teil 5 Kinder setzen sich auf 5 Stühle in einer Reihe. Wie viele Sitzordnungen gibt es?
\teil Bilde aus den Ziffern 3, 9 und 5 – jede genau einmal – die größte und die kleinste dreistellige Zahl.
\teil Zwei gleiche Spielsteine mit einer roten und einer blauen Seite werden gleichzeitig geworfen. Wie viele verschiedene Ergebnisse sind möglich? \hfill (P10 2017 OS)
\teil Ein Würfel wird zweimal geworfen. Schreibe alle Paare auf, bei denen der zweite Wurf eine 4 zeigt. \hfill (P10 2020 OS)
\teil Die vier Punkte A, B, C, D liegen so, dass A, B und C auf einer Geraden liegen. Wie viele Dreiecke mit drei dieser Punkte als Ecken gibt es?
\teil Ein Fahrradschloss hat einen dreistelligen Code aus den Ziffern 2, 4 und 8; jede Ziffer kommt genau einmal vor. Wie viele Codes sind möglich? \hfill (P10 2017 OS)
\end{teile}
\end{aufgabe}

%% TODO Titel: Ich-kann-Satz fehlt in der Bank (Platzhalter: Kettenname) – Nr. 9
%% TODO Anweisung: ein Satz über der ersten Teilaufgabe fehlt in der Bank – Nr. 9
\begin{aufgabe}{Ergebnisse zu „mindestens einmal …“ aufzählen}
\begin{teile}
\teil Eine Münze wird zweimal geworfen. Schreibe alle Ergebnisse auf, bei denen mindestens einmal Kopf fällt.
\end{teile}
\end{aufgabe}

%% TODO Titel: Ich-kann-Satz fehlt in der Bank (Platzhalter: Kettenname) – Nr. 10
%% TODO Anweisung: ein Satz über der ersten Teilaufgabe fehlt in der Bank – Nr. 10
\begin{aufgabe}{Zählen – Fehler finden, Begründen, Darstellungswechsel, Anwendung}
\begin{teile}
\teil Ole wirft zwei Münzen und schreibt: „Mögliche Ergebnisse sind (Kopf; Kopf), (Kopf; Zahl) und (Zahl; Zahl), also 3.“ Finde den Fehler.
\teil Sara schreibt alle Eisbecher aus einer Sorte (Zitrone, Kirsch) und einer Waffel (klein, groß) geordnet auf: Zitrone–klein, Zitrone–groß, Kirsch–klein, Kirsch–groß. Begründe, warum keine Möglichkeit fehlt.
\teil Der Baum zeigt die Wahl eines Getränks und eines Kuchens. Schreibe alle Möglichkeiten als Liste auf.

\baumzwei{Saft/,Tee/}{Apfel/,Nuss/}{Apfel/,Nuss/}
\teil Ein Kofferschloss hat drei Ringe mit je den Ziffern 0 bis 9. Du probierst jede Sekunde einen Code. Wie lange dauert es höchstens, bis du alle Codes durchprobiert hast?
\end{teile}
\end{aufgabe}
